% part: intuitionistic-logic
% chapter: soundness-completeness
% section: decidability

\documentclass[../../../include/open-logic-section]{subfiles}

\begin{document}

\olfileid{int}{sc}{dec}

\olsection{القابلية للقرار}

كلما كان~$\OLId{greek-variable}{Gamma} \Proves/ !A$ أعطانا برهان التمام نموذجًا عوالمه غير
متناهية، يشهد بـ~$\OLId{greek-variable}{Gamma} \Entails/ !A$. أما النتيجة الآتية فتبيّن أن
إقامة~$\Entails !A$ لا تحتاج إلا إلى إقامة
$\mSat{\OLId{latin-looped}{M}}{!A}$ في النماذج المنتهية جميعًا؛ ونعني بالمنتهي ما كانت
مجموعة عوالمه منتهية.

\begin{thm}\ollabel{thm:decidability}
  متى كان~$\Entails/ !A$ وُجد نموذج منته~$\mModel{M'}$ يحقق
  $\mSat/{\OLId{latin-looped}{M}'}{!A}$.
\end{thm}

\begin{proof}
  ليكن~$\mModel{M}=\tuple{\OLId{latin-looped}{W},\OLId{latin-looped}{R},\OLId{latin-looped}{V}}$ نموذجًا، ولنأخذ~$\OLId{latin-ordinary}{w}_\OLMathNumeral{0}\in \OLId{latin-looped}{W}$ بحيث
  $\mSat/{\OLId{latin-looped}{M}}{!A}[\OLId{latin-ordinary}{w}_\OLMathNumeral{0}]$. نسمّي~$\OLId{latin-looped}{P}$ مجموعة ال!!{propositional variable}s
  الواقعة في~$!A$، ونسمي~$\Sigma$ إغلاق~$\{!A\}$ بأخذ
  الصيغ الجزئية: فـ~$\Sigma$ تضم الصيغ الجزئية
  لـ~$!A$ كلها، وهي لذلك منتهية.

  ونحدّ علاقة التكافؤ~$\sim_\Sigma$ على~$\OLId{latin-looped}{W}$ كما يأتي:
  \[
    \OLId{latin-ordinary}{u}\sim_\Sigma \OLId{latin-ordinary}{v}
    \quad\text{إذا وفقط إذا}\quad
    \bigl(\mSat{\OLId{latin-looped}{M}}{!B}[\OLId{latin-ordinary}{u}]\ \text{إذا وفقط إذا}\
    \mSat{\OLId{latin-looped}{M}}{!B}[\OLId{latin-ordinary}{v}]\bigr)
    \quad\text{لكل }!B\in\Sigma.
  \]
  ولتكن~$[\OLId{latin-ordinary}{u}]_\Sigma$ فئة تكافؤ~$\OLId{latin-ordinary}{u}$. ثم نعرّف
  $\mModel{M'}=\tuple{\OLId{latin-looped}{W}',\OLId{latin-looped}{R}',\OLId{latin-looped}{V}'}$ بهذه الحدود:
  \[
    \OLId{latin-looped}{W}'=\Setabs{[\OLId{latin-ordinary}{u}]_\Sigma}{\OLId{latin-ordinary}{u}\in \OLId{latin-looped}{W}},
  \]
  \[
    \OLId{latin-looped}{R}'{[\OLId{latin-ordinary}{u}]_\Sigma}{[\OLId{latin-ordinary}{v}]_\Sigma}
    \quad\text{إذا وفقط إذا}\quad
    \bigl(\mSat{\OLId{latin-looped}{M}}{!B}[\OLId{latin-ordinary}{u}]\ \text{يستلزم}\
    \mSat{\OLId{latin-looped}{M}}{!B}[\OLId{latin-ordinary}{v}]\bigr)
    \quad\text{لكل }!B\in\Sigma,
  \]
  و
  \[
    \OLId{latin-looped}{V}'(\OLId{latin-ordinary}{p})=
    \begin{cases}
      \Setabs{[\OLId{latin-ordinary}{u}]_\Sigma}{\OLId{latin-ordinary}{u}\in \OLId{latin-looped}{W}\text{ و }\mSat{\OLId{latin-looped}{M}}{\OLId{latin-ordinary}{p}}[\OLId{latin-ordinary}{u}]}
        & \text{إذا كان }\OLId{latin-ordinary}{p}\in\Sigma,\\
      \varnothing & \text{خلاف ذلك.}
    \end{cases}
  \]
  وليس لاختيار ممثل الفئة أثر في هذه التعريفات. وأما~$\OLId{latin-looped}{R}'$
  فانعكاسيتها وتعدّيها ظاهران. فإذا اجتمع
  $\OLId{latin-looped}{R}'{[\OLId{latin-ordinary}{u}]_\Sigma}{[\OLId{latin-ordinary}{v}]_\Sigma}$ و
  $\OLId{latin-looped}{R}'{[\OLId{latin-ordinary}{v}]_\Sigma}{[\OLId{latin-ordinary}{u}]_\Sigma}$، اتفق~$\OLId{latin-ordinary}{u}$ و$\OLId{latin-ordinary}{v}$ في صدق الصيغ كلها
  من~$\Sigma$، فيكون~$[\OLId{latin-ordinary}{u}]_\Sigma=[\OLId{latin-ordinary}{v}]_\Sigma$. فثبت أن~$\OLId{latin-looped}{R}'$ ترتيب
  جزئي. و~$\OLId{latin-looped}{V}'$ رتيبة أيضًا بالنسبة إلى~$\OLId{latin-looped}{R}'$؛ فإذن
  $\mModel{M'}$ هو !!a{relational model}.

  وكل فئة في~$\OLId{latin-looped}{W}'$ يعيّنها ما يصدق من صيغ~$\Sigma$ عند ممثليها؛
  ومن ذلك نحصل على الحد
  \[
    |\OLId{latin-looped}{W}'|\leq \OLMathNumeral{2}^{|\Sigma|}.
  \]
  فـ~$\OLId{latin-looped}{W}'$ منتهية على الخصوص.

  ونقيم الآن بالاستقراء البنيوي لمّة الترشيح لكل~!!{formula}~$!B$
  في~$\Sigma$:
  \[
    \mSat{\OLId{latin-looped}{M}}{!B}[\OLId{latin-ordinary}{u}]
    \quad\text{إذا وفقط إذا}\quad
    \mSat{\OLId{latin-looped}{M}'}{!B}[{[\OLId{latin-ordinary}{u}]_\Sigma}]
    \qquad(!B\in\Sigma).
  \]
  أما ال!!{propositional variable}s و$\lfalse$ و$\land$ و$\lor$
  فحالاتها تأتي بلا واسطة
  من تعريف~$\OLId{latin-looped}{V}'$ وفرض الاستقراء. والمحتاج إلى بيان هو حفظ النفي والشرط
  عند هذا الاختيار لـ~$\OLId{latin-looped}{R}'$.

  لنأخذ حالة~$!B\ident !C\lif !D$. إذا كان أولًا
  $\mSat{\OLId{latin-looped}{M}}{!C\lif !D}[\OLId{latin-ordinary}{u}]$، وأخذنا
  $\OLId{latin-looped}{R}'{[\OLId{latin-ordinary}{u}]_\Sigma}{[\OLId{latin-ordinary}{v}]_\Sigma}$، فإن
  $!C\lif !D\in\Sigma$؛ وتعريف~$\OLId{latin-looped}{R}'$ ينقل صدقها إلى~$\OLId{latin-ordinary}{v}$،
  فيثبت~$\mSat{\OLId{latin-looped}{M}}{!C\lif !D}[\OLId{latin-ordinary}{v}]$. فإذا تحقق
  $\mSat{\OLId{latin-looped}{M}'}{!C}[{[\OLId{latin-ordinary}{v}]_\Sigma}]$، ثبت بالاستقراء
  $\mSat{\OLId{latin-looped}{M}}{!C}[\OLId{latin-ordinary}{v}]$؛ وانعكاسية~$\OLId{latin-looped}{R}$ تعطينا حينئذ
  $\mSat{\OLId{latin-looped}{M}}{!D}[\OLId{latin-ordinary}{v}]$. ثم نطبق فرض الاستقراء ثانيةً فنحصل على
  $\mSat{\OLId{latin-looped}{M}'}{!D}[{[\OLId{latin-ordinary}{v}]_\Sigma}]$. وبهذا يثبت
  $\mSat{\OLId{latin-looped}{M}'}{!C\lif !D}[{[\OLId{latin-ordinary}{u}]_\Sigma}]$.

  وأما العكس، فلو كان~$\mSat/{\OLId{latin-looped}{M}}{!C\lif !D}[\OLId{latin-ordinary}{u}]$ لكان عالم~$\OLId{latin-ordinary}{v}$ يحقق
  $Ruv$ و$\mSat{\OLId{latin-looped}{M}}{!C}[\OLId{latin-ordinary}{v}]$ و$\mSat/{\OLId{latin-looped}{M}}{!D}[\OLId{latin-ordinary}{v}]$. ورتابة الصدق في
  النموذج الأول تقتضي انتقال الصدق من~$\OLId{latin-ordinary}{u}$ إلى~$\OLId{latin-ordinary}{v}$ لكل صيغة من~$\Sigma$،
  فيثبت~$\OLId{latin-looped}{R}'{[\OLId{latin-ordinary}{u}]_\Sigma}{[\OLId{latin-ordinary}{v}]_\Sigma}$. ويعطي فرض الاستقراء صدق~$!C$
  وعدم صدق~$!D$ عند~$[\OLId{latin-ordinary}{v}]_\Sigma$ في~$\mModel{M'}$؛ فينتج
  $\mSat/{\OLId{latin-looped}{M}'}{!C\lif !D}[{[\OLId{latin-ordinary}{u}]_\Sigma}]$.

  وأما حالة~$!B\ident\lnot !C$ فوجهها نظير ذلك. فمن
  $\mSat{\OLId{latin-looped}{M}}{\lnot !C}[\OLId{latin-ordinary}{u}]$ و
  $\OLId{latin-looped}{R}'{[\OLId{latin-ordinary}{u}]_\Sigma}{[\OLId{latin-ordinary}{v}]_\Sigma}$ ينقل تعريف~$\OLId{latin-looped}{R}'$ صدق
  $\lnot !C\in\Sigma$ إلى~$\OLId{latin-ordinary}{v}$. ولانعكاسية~$\OLId{latin-looped}{R}$ لا يصدق~$!C$ عند~$\OLId{latin-ordinary}{v}$؛
  فلا يصدق بالاستقراء عند~$[\OLId{latin-ordinary}{v}]_\Sigma$ في~$\mModel{M'}$. وإن
  كان~$\mSat/{\OLId{latin-looped}{M}}{\lnot !C}[\OLId{latin-ordinary}{u}]$، وُجد~$\OLId{latin-ordinary}{v}$ يحقق~$Ruv$ ويصدق عنده~$!C$
  أي عند~$\OLId{latin-ordinary}{v}$. فرتابة الصدق تعطينا~$\OLId{latin-looped}{R}'{[\OLId{latin-ordinary}{u}]_\Sigma}{[\OLId{latin-ordinary}{v}]_\Sigma}$،
  ويعطينا فرض الاستقراء شاهدًا يثبت
  $\mSat/{\OLId{latin-looped}{M}'}{\lnot !C}[{[\OLId{latin-ordinary}{u}]_\Sigma}]$.

  وقد كانت ال!!{formula}~$!A$ من~$\Sigma$؛ فتنتج لمّة الترشيح
  $\mSat/{\OLId{latin-looped}{M}'}{!A}[{[\OLId{latin-ordinary}{w}_\OLMathNumeral{0}]_\Sigma}]$. فصار~$\mModel{M'}$ نموذجًا
  مضادًا منتهيًا لـ~$!A$، وتم المطلوب.
\end{proof}

\begin{prob}
أتمّ التفاصيل المباشرة في برهان~\olref[int][sc][dec]{thm:decidability}:
أثبت حسن تعريف~$\OLId{latin-looped}{R}'$ و$\OLId{latin-looped}{V}'$، وكون~$\OLId{latin-looped}{R}'$ ترتيبًا جزئيًا و~$\OLId{latin-looped}{V}'$
رتيبة، ثم أقم حالات المتغيرات القضوية و$\lfalse$ و$\land$ و$\lor$ في
لمّة الترشيح.
\end{prob}

وتعطي~\olref{thm:decidability} طريقًا خوارزميًا للفصل في
$\Entails !A$: نقتصر في الفحص على النماذج العلائقية التي عدد عوالمها
لا يزيد على~$\OLMathNumeral{2}^{|\Sigma|}$، وعلى المتغيرات القضوية الواقعة
في~$!A$.

\end{document}
